\documentclass[11pt,a4paper]{article}
\usepackage{../pelm}
\begin{document}
\paper{Quiz 4}{Weeks 7--8: Grounding, and Verification}{10}{10 minutes}

\begin{questions}

\q{``Grounding'' a system in your own sources means:}
  {Retraining the model on your documents}
  {Placing relevant passages from your documents into the context when it answers}
  {Restricting the model to a single topic}
  {Reducing the model's temperature}

\q{The abbreviation RAG stands for retrieval-augmented generation. In practice it means:}
  {A method for training smaller models}
  {Retrieve relevant text, add it to the context, then generate}
  {A technique for reducing token costs}
  {A way of checking output against the internet}

\q{Which problem does grounding \textbf{not} solve?}
  {The model has never seen your institution's internal policy}
  {The guideline was updated after the knowledge cut-off}
  {Your source document is itself out of date}
  {Turkish-language professional material is thinly represented}

\q{Why should the refusal instruction specify a fixed string such as ``NOT IN SOURCES''?}
  {It reduces the number of tokens used}
  {It is easy to spot when skim-reading, unlike a hedge buried in a paragraph}
  {It prevents the model from answering at all}
  {It is required by most document tools}

\q{An assistant answers ``The notice period is 30 days'' with no quotation from your documents.
   Which grounded failure mode should you suspect first?}
  {Over-read the passage}
  {Retrieved the wrong passage}
  {Answered from memory}
  {Silent failure}

\q{Your source states: ``Either party \emph{may} request an extension.'' The assistant reports:
   ``Either party \emph{must} request an extension.'' This is:}
  {Fabrication}
  {Over-reading the passage}
  {Answered from memory}
  {Displacement}

\q{In a ten-question test set, what is the recommended composition?}
  {Ten questions your sources answer}
  {Seven your sources answer and three they cannot}
  {Five easy and five hard}
  {Ten questions chosen at random}

\q{Why do the unanswerable questions matter more than the answerable ones?}
  {They are harder to write}
  {They test behaviour in the situation you cannot detect during normal use}
  {They take less time to grade}
  {They are more likely to appear in the exam}

\q{Why must you write the expected answer \textbf{before} running each test question?}
  {It saves time}
  {Otherwise you will unconsciously accept whatever output you receive}
  {The tool requires it}
  {It improves retrieval accuracy}

\q{Which of these is genuine verification?}
  {Asking the same question again and getting a similar answer}
  {Asking the model whether it is confident}
  {Searching your source document for the quoted sentence}
  {Regenerating the answer three times and taking the majority}

\end{questions}
\end{document}
